% Shared preamble for all outputs (master + standalone targets)
\usepackage[a4paper,twoside,margin=22mm,includeheadfoot]{geometry}
\usepackage[T1]{fontenc}
\usepackage[utf8]{inputenc}
\usepackage{textcomp}
\usepackage{lmodern}

\usepackage{microtype}
\usepackage{hyperref}
\hypersetup{hidelinks}

\usepackage{xcolor}
\usepackage{graphicx}
\usepackage{enumitem}
\usepackage{booktabs}
\usepackage{tabularx}
\usepackage{array}
\usepackage{longtable}
\usepackage{multicol}
\usepackage{etoolbox}

\usepackage{titlesec}
\usepackage{fancyhdr}
\usepackage{lastpage}
\usepackage[absolute,overlay]{textpos}

% --- Minimal man-page-ish styling ---
\setlength{\parindent}{0pt}
\setlength{\parskip}{0.35em}

\titleformat{\chapter}[display]
  {\normalfont\bfseries\Large}
  {\MakeUppercase{Book \thechapter}}
  {0.3em}
  {\Huge}
\titlespacing*{\chapter}{0pt}{0.8em}{0.6em}

\titleformat{\section}
  {\normalfont\bfseries\Large}
  {\thesection}
  {0.6em}
  {}
\titleformat{\subsection}
  {\normalfont\bfseries\normalsize}
  {\thesubsection}
  {0.6em}
  {}
\titleformat{\subsubsection}
  {\normalfont\bfseries\normalsize}
  {\thesubsubsection}
  {0.6em}
  {}

% No printed page numbers (binder replacement workflow)
\pagestyle{fancy}
\fancyhf{}
\renewcommand{\headrulewidth}{0.4pt}
\renewcommand{\footrulewidth}{0.4pt}

% Header shows navigation context; footer shows version + credits.
\fancyhead[LE]{\small\ttfamily \leftmark}
\fancyhead[RO]{\small\ttfamily \rightmark}
\fancyhead[LO]{\small\ttfamily TAB: \CurrentTabLabel}
\fancyhead[RE]{\small\ttfamily TAB: \CurrentTabLabel}

\fancyfoot[LE,RO]{\small\ttfamily v\MatterVersion\; (\MatterDate)}
\fancyfoot[LO,RE]{\small\ttfamily credit: \CurrentCredit}

% Chapter opening pages use plain style - make it empty (no page numbers)
\fancypagestyle{plain}{%
  \fancyhf{}%
  \renewcommand{\headrulewidth}{0pt}%
  \renewcommand{\footrulewidth}{0pt}%
}

% Keep headings from forcing lowercase marks
\makeatletter
\@ifundefined{chapter}{%
  % article-like classes: only section-level marks exist
  \renewcommand{\sectionmark}[1]{\markboth{#1}{#1}}%
}{%
  % book/report classes: keep chapter + section context
  \renewcommand{\chaptermark}[1]{\markboth{#1}{#1}}%
  \renewcommand{\sectionmark}[1]{\markright{#1}}%
}
\makeatother

% --- Title page styling macros (formal manual style) ---
\newcommand{\TitleLarge}[1]{{\Huge\bfseries\rmfamily #1}}
\newcommand{\TitleMedium}[1]{{\Large\bfseries\rmfamily #1}}
\newcommand{\TitleSmall}[1]{{\normalsize\rmfamily #1}}
\newcommand{\TitleMono}[1]{{\normalsize\ttfamily #1}}
\newcommand{\TitleThin}{\rule{0.6\textwidth}{0.4pt}}
\newcommand{\TitleRule}{\rule{0.5\textwidth}{0.6pt}}

% --- Margin tab marker ---
% A restrained default palette; can be overridden centrally.
\definecolor{TabGray}{gray}{0.85}

\newcommand{\tabmarker}[1]{%
  \marginpar[\raggedleft\footnotesize\ttfamily\colorbox{TabGray}{\strut\,#1\,}]
            {\raggedright\footnotesize\ttfamily\colorbox{TabGray}{\strut\,#1\,}}%
}

% --- Hierarchical tab-aware navigation list ---
\newcommand{\MatterNavEntries}{}
\newcommand{\MatterNavLine}[3]{%
  \par\noindent
  \ifcase#1\relax
    \textbf{#2}\dotfill\texttt{#3}%
  \or
    \hspace*{1.2em}#2\dotfill\texttt{#3}%
  \or
    \hspace*{2.4em}#2\dotfill\texttt{#3}%
  \or
    \hspace*{3.6em}#2\dotfill\texttt{#3}%
  \else
    #2\dotfill\texttt{#3}%
  \fi
  \par
}

% --- Book registry for title page (NATOPS-style chapter listing) ---
\newcommand{\MatterBookEntries}{}
\newcommand{\MatterBookBoxEntries}{}
\newcommand{\MatterBookLine}[2]{%
  \par\noindent
  #1\dotfill\texttt{#2}\par
}
\newcommand{\MatterBookLineAsBox}[2]{%
  \colorbox[HTML]{353535}{\makebox[35mm][c]{%
    \parbox[c][20mm]{33mm}{\centering
      \vspace{0.4em}
      \textcolor{white}{\fontsize{11pt}{13pt}\selectfont #1}\\[0.3cm]
      \textcolor{white}{\fontsize{10pt}{11pt}\selectfont \texttt{#2}}
      \vspace{0.4em}
    }%
  }}\par\vspace{0.8em}%
}

\makeatletter
\newcommand{\RecordMatterBook}[2]{%
  \g@addto@macro\MatterBookEntries{\MatterBookLine{#1}{#2}}%
  \g@addto@macro\MatterBookBoxEntries{\MatterBookLineAsBox{#1}{#2}}%
}
\newcommand{\MatterBookDirect}[2]{%
  \protected@write\@auxout{}{%
    \string\RecordMatterBook{#1}{#2}%
  }%
}
\newcommand{\AddMatterNavEntry}[3]{\g@addto@macro\MatterNavEntries{\MatterNavLine{#1}{#2}{#3}}}
\newcommand{\RecordMatterNavEntry}[2]{%
  \AddMatterNavEntry{#1}{#2}{\CurrentTabLabel}%
  \protected@write\@auxout{}{%
    \string\AddMatterNavEntry{#1}{#2}{\CurrentTabLabel}%
  }%
}
\makeatother

\newcommand{\listofbooks}{%
  \ifstrempty{\MatterBookEntries}{}{%
    \MatterBookEntries
  }%
}

\newcommand{\listofbooksasboxes}{%
  \ifstrempty{\MatterBookBoxEntries}{}{%
    \vfill
    \MatterBookBoxEntries
    \vfill
  }%
}

% --- Per-book table of contents ---
\newcommand{\CurrentBookTOCEntries}{}
\newcommand{\CurrentBookLabel}{}

% Global TOC registry for each book (persists across builds)
\newcommand{\BookTOCFND}{}

% Simple counter-based TOC line renderer (avoids nested enumerate complexity)
\newcounter{bTOClvlI}
\newcounter{bTOClvlII}
\newcounter{bTOClvlIII}
\newcommand{\BookTOCLine}[3]{%
  % #1 = level (0-3), #2 = title, #3 = label (unused)
  \ifcase#1\relax
    % Level 0: do nothing
  \or
    % Level 1
    \stepcounter{bTOClvlI}%
    \setcounter{bTOClvlII}{0}%
    \setcounter{bTOClvlIII}{0}%
    \par\noindent\textbf{\arabic{bTOClvlI}.}\quad #2\par
  \or
    % Level 2
    \stepcounter{bTOClvlII}%
    \setcounter{bTOClvlIII}{0}%
    \par\noindent\hspace*{1.2em}\textbf{\alph{bTOClvlII}.}\quad #2\par
  \or
    % Level 3
    \stepcounter{bTOClvlIII}%
    \par\noindent\hspace*{2.4em}\textbf{\roman{bTOClvlIII}.}\quad #2\par
  \fi
}

\makeatletter
\newcommand{\AddBookTOCEntry}[4]{%
  % #1 = book label, #2 = level, #3 = title, #4 = unused
  % Append to the specified book's TOC macro
  \expandafter\g@addto@macro\csname BookTOC#1\endcsname{\BookTOCLine{#2}{#3}{#4}}%
}
\newcommand{\RecordBookTOCEntry}[3]{%
  \AddBookTOCEntry{\CurrentBookLabel}{#1}{#2}{#3}%
  \protected@write\@auxout{}{%
    \string\AddBookTOCEntry{\CurrentBookLabel}{#1}{#2}{#3}%
  }%
}
\newcommand{\PrintBookTOC}{%
  \expandafter\ifx\csname BookTOC\CurrentBookLabel\endcsname\relax
    % Macro doesn't exist
  \else
    \ifstrempty{\csname BookTOC\CurrentBookLabel\endcsname}{}{%
      \section*{Contents}%
      \small
      % Execute the custom TOC lines in standard text mode
      \csname BookTOC\CurrentBookLabel\endcsname
      \normalsize
    }%
  \fi
}
\makeatother

\newcommand{\listoftabs}{%
  \section*{Navigation}
  \begingroup
    \parskip=0.2em
    \ifstrempty{\MatterNavEntries}{\noindent\emph{(Navigation will appear after next run.)}\par}{\MatterNavEntries}
  \endgroup
}

% Heading wrappers so the navigation list is complete and tab-aware.
\newcommand{\MatterBook}[2]{%
  \settablabel{#2}%
  \renewcommand{\CurrentBookLabel}{#2}%
  \MatterBookDirect{#1}{#2}%
  \chapter{#1}%
  \PrintBookTOC%
  \RecordMatterNavEntry{0}{#1}%
}

\newcommand{\MatterChapter}[1]{%
  \section{#1}%
  \RecordBookTOCEntry{1}{#1}{}%
  \RecordMatterNavEntry{1}{#1}%
}

\newcommand{\MatterSection}[1]{%
  \subsection{#1}%
  \RecordBookTOCEntry{2}{#1}{}%
  \RecordMatterNavEntry{2}{#1}%
}

% --- Replacement-friendly hierarchy helpers ---
% Subsubsection is treated as the 'Subsection' level for replacement.
% Each Subsection starts on a fresh page so it can be reprinted independently.
\newcommand{\MatterSubsection}[3]{%
  % #1 = label, #2 = title, #3 = credit
  \clearpage
  \setcredit{#3}
  \subsubsection{#2}\label{#1}
  \RecordBookTOCEntry{3}{#2}{}%
  \RecordMatterNavEntry{3}{#2}%
  \tabmarker{\CurrentTabLabel}
}

\makeatletter
% 1. The Token-Splitting Engine
\def\FakeSC#1{\@FakeSC#1\@nil}
\def\@FakeSC#1#2\@nil{\textbf{\MakeUppercase{#1}{\scriptsize\MakeUppercase{#2}}}}
\makeatother

% 2. The Custom Sectioning Wrapper
\newcommand{\msubsubsection}[1]{%
  % The optional argument [#1] ensures the TOC gets plain text.
  % The required argument {\FakeSC{#1}} prints the stylised text.
  \subsubsection[#1]{\FakeSC{#1}}%
}

% --- Checklist Environment & Row Command ---
% --- Checklist Environment & Row Command ---

\newenvironment{checklist}[1]{%
  % 1. Print the title
  \par\vspace*{1em}\noindent\textbf{#1}\par\nopagebreak\vspace*{0.5em}
  % 2. Open the table using the primitive \tabularx
  \noindent
  \renewcommand{\arraystretch}{1.5}
  \tabularx{\textwidth}{|>{\raggedright\arraybackslash}p{0.22\textwidth}|
                                 >{\raggedright\arraybackslash}X|
                                 >{\raggedright\arraybackslash}X|}
  \hline
  % 3. The locked-in headers
  \multicolumn{1}{|c|}{\textbf{ITEM}} & 
  \multicolumn{1}{c|}{\textbf{DESCRIPTION}} & 
  \multicolumn{1}{c|}{\textbf{EXAMPLE}} \\
  \hline
}{%
  % 4. Close the table using the primitive \endtabularx
  \endtabularx\par\vspace*{1em}
}

% The data entry command remains unchanged
\newcommand{\checkrow}[3]{%
  #1 & #2 & #3 \\ \hline
}

% Lists: compact and readable
\setlist[itemize]{leftmargin=*,itemsep=0.15em,topsep=0.25em}
\setlist[enumerate]{leftmargin=*,itemsep=0.15em,topsep=0.25em}
